\chapter{Parabolic and hyperbolic problems: transient analysis}
\label{chap:heat}

\framepart{Outline}

This chapter adds time. We take the heat conduction equations in the
transient regime and study the heating of the plate with a hole, then do the
same for elastodynamics with the Newmark algorithm.

Both problems could be handed to the \op{pasapas} operator, which is what you
would do in production and what chapter~\ref{chap:plasticity} describes. We
program the time loop explicitly here instead, because writing the algorithm
out in \textsc{gibiane} is the quickest way to understand what \op{pasapas}
does for you -- and because the same loop structure is what you will reuse
for any algorithm of your own.

\section{Transient heat transfer}

\noindent This section continues the thermal problem of
section~\ref{sec:elas:thermal}: the model \op{mod}, the material \op{mat}
and the blocked boundary \op{Ad23} are the ones built there, so run that
section first (or restore its objects with \op{rest}).


Recall that the family of algorithms presented in the lecture notes rests on
the equation

\begin{equation}
\left[\frac{1}{\Delta t}M + \theta K\right] T_{n+1}
=
\left[\frac{1}{\Delta t}M - (1-\theta) K\right] T_{n}
+ \theta F_{n+1} + (1-\theta) F_{n}
\label{eq:heat:theta}
\end{equation}

\cindex{transient analysis}\cindex{time integration}%
\cindex[theta-method]{$\theta$-method}\cindex{Crank--Nicolson scheme}%
\cindex{implicit scheme}\cindex{explicit scheme}\cindex{stability, unconditional}%
\noindent which determines $T_{n+1}$ knowing $T_{n}$. Here $M$ is the
capacity matrix, $K$ the conductivity matrix, and $\theta\in[0,1]$ is a
parameter: $\theta=0$ gives the explicit Euler scheme, $\theta=1$ the implicit
one, $\theta=1/2$ the Crank--Nicolson scheme, which is second-order accurate.
Only $\theta\geq1/2$ is unconditionally stable.

\begin{castemcom}
We start by fixing the time step and the algorithm parameter. The choice
$\theta=1$ gives an implicit and therefore unconditionally stable algorithm --
the safe default while you are debugging.

Note that with $\theta=1$ there is no restriction on $\Delta t$ for stability,
but there very much is one for accuracy: the solution will be smooth and
wrong rather than oscillatory and wrong.
\end{castemcom}
\begin{castemlines}
\begin{verbatim}
dt    = 0.1;
theta = 1.;
\end{verbatim}
\end{castemlines}

\begin{castemcom}
\sindex{table, as a history container}\cindex{time step}%
The conductivity and capacity matrices are built first. Both are objects of
\op{rigidite} type, which is why they can be combined with \op{et} and
multiplied by scalars exactly like stiffness matrices.

The two matrices of the algorithm follow directly from
\eqref{eq:heat:theta}. They are assembled once, outside the loop -- this is
the whole point of a constant time step, since \op{Matnpun} is the only
matrix that has to be factorised.
\end{castemcom}
\begin{castemlines}
\begin{verbatim}
KK = cond mod mat;
MM = capa mod mat;

Matnpun = ((1./dt) * MM)
          et (theta * KK);
Matn    = ((1./dt) * MM)
          et ((theta - 1.) * KK);
\end{verbatim}
\end{castemlines}

\begin{castemcom}
The matrix $A$ corresponding to the temperature imposed on the exterior
boundary is built as in the steady case, and the maximum exterior temperature
\op{Td} is fixed.
\end{castemcom}
\begin{castemlines}
\begin{verbatim}
Td = 1.;

Ad23 = bloq 'T' d23;
AA   = Ad23;
\end{verbatim}
\end{castemlines}

\begin{castemcom}
The tables of temperature fields and of imposed temperatures are initialised.
A \op{table} indexed by the step number is the natural \textsc{gibiane}
container for a history: \op{T . n} is the field at step $n$, and the whole
history stays available for post-processing after the loop.
\end{castemcom}
\begin{castemlines}
\begin{verbatim}
T = table;
T . 0 = manu 'CHPO' dom 1 'T' 0.;

F = table;
F . 0 = depi Ad23 0.;
\end{verbatim}
\end{castemlines}

\begin{castemcom}
The computation is then the resolution of \eqref{eq:heat:theta} at each time
step.

Two points of \textsc{gibiane} syntax are worth noting. \op{\&boucle} is the
iterator of the loop named \op{boucle}, and it runs from 1, which is why
\sindex{\&, loop iterator}%
\op{n} is defined as \op{\&boucle - 1} so that the indices match the tables
starting at 0. And the index of a table may be any expression, provided it is
bracketed: \op{F . (n + 1)}.

Here the imposed boundary temperature ramps linearly from 0 to \op{Td} over
the \op{ndt} steps.

Note the two different combinations in the right-hand side: \op{+} adds two
nodal fields defined on the same degrees of freedom, whereas \op{et}
\emph{concatenates} contributions living on different ones -- here the
volume term and the Lagrange-multiplier term of the blocked boundary. They
are not interchangeable.
\end{castemcom}
\begin{castemlines}
\begin{verbatim}
ndt = 20;

repeter boucle ndt;
  n = &boucle - 1;

  F . (n + 1) = depi Ad23
                ((n + 1) * Td / ndt);

  Fn = (theta * (F . (n + 1))) +
       ((1. - theta) * (F . n));

  T . (n + 1) = reso (Matnpun et AA)
                ((Matn * (T . n)) et Fn);
fin boucle;
\end{verbatim}
\end{castemlines}

\begin{castemcom}
After the loop the history is in the table, and any point value can be
followed in time by extracting it step by step into a \op{listreel} and
building an \op{evolution}.

\op{extr} with a point as last argument takes the value of a nodal field at
that point. This is the idiom you will use over and over for time histories.
\end{castemcom}
\begin{castemlines}
\begin{verbatim}
tt   = prog;
thi  = prog;
pA   = dom poin proc (0.5 0.5);

repeter post (ndt + 1);
  n = &post - 1;
  tt  = tt  et (prog (n * dt));
  thi = thi et
     (prog (extr (T . n) 'T' pA));
fin post;

dess (evol manu 'time' tt
               'T' thi);
\end{verbatim}
\end{castemlines}

\section{Example in dynamic analysis}

\cindex{dynamics, elastic}\cindex{Newmark algorithm}%
\cindex{average acceleration scheme}\sindex{dyne}%
\cindex{modal projection}\cindex{wave propagation}%
As an example of dynamic analysis we take the equations of linear elasticity
in the dynamic regime and program the Newmark algorithm in \textsc{gibiane}.

As for the transient thermal problem, this type of problem can be solved with
the \op{pasapas} operator. A variant based on the projection of the
displacement field on a modal basis is available through the \op{dyne}
operator, which is far cheaper when many time steps are needed and the
response is dominated by a few modes.

The Newmark family rests on the prediction/correction pair

\begin{align}
\tilde{U}_{n+1} &= U_{n} + \Delta t\,\dot{U}_{n}
   + \tfrac{1}{2}\Delta t^{2}(1-2\beta)\,\ddot{U}_{n} \\
\dot{\tilde{U}}_{n+1} &= \dot{U}_{n} + \Delta t (1-\gamma)\,\ddot{U}_{n}
\end{align}

\noindent followed by the resolution of
$\left[M+\beta\Delta t^{2}K\right]\ddot{U}_{n+1}
 = F_{n+1}-K\tilde{U}_{n+1}$
and the update
$U_{n+1}=\tilde{U}_{n+1}+\beta\Delta t^{2}\ddot{U}_{n+1}$,
$\dot{U}_{n+1}=\dot{\tilde{U}}_{n+1}+\gamma\Delta t\,\ddot{U}_{n+1}$.

\begin{castemcom}
A bar of length \op{ldom} and height \op{hdom} is built starting from the line
defining its left boundary, \op{dgauche}, and generating the mesh with
\op{tran[slation]}. The right boundary is recovered as face 3 of the
generated surface.
\end{castemcom}
\begin{castemlines}
\begin{verbatim}
* beware: this redefines ldom, which
* chapter 3 also uses
ldom = 50.; hdom = 1.;
p1 = 0. 0.; p2 = 0. hdom;

nlelem = 100; nhelem = 2;

dgauche = droi nhelem p1 p2;
dom     = dgauche tran nlelem
                  (ldom 0.);
ddroite = dom face 3;
\end{verbatim}
\end{castemlines}

\begin{castemcom}
The model and the material follow. Note that \texttt{'RHO'} is compulsory
here: without a density there is no mass matrix and no dynamics.

With $E=1$, $\rho=1$ the wave speed is $c=\sqrt{E/\rho}=1$, so a wave crosses
the 50-unit bar in 50 time units -- convenient for checking the results
against the analytical solution.
\end{castemcom}
\begin{castemlines}
\begin{verbatim}
mod = dom mode mecanique
          elastique isotrope;
mat = mod mate 'RHO' 1.
               'YOUN' 1. 'NU' 0.3;
\end{verbatim}
\end{castemlines}

\begin{castemcom}
The parameters of the computation: the time step \op{dt} and the algorithm
parameters $\beta$ and $\gamma$.

The values $\beta=1/4$, $\gamma=1/2$ give the average acceleration scheme,
which is implicit, second-order accurate, unconditionally stable and
conserves energy. They are the standard choice and the one to use unless you
have a reason not to.
\end{castemcom}
\begin{castemlines}
\begin{verbatim}
ndt   = 200;
dt    = 1.;
beta  = 0.25;
gamma = 0.5;
\end{verbatim}
\end{castemlines}

\begin{castemcom}
The boundary conditions: the left end of the bar is clamped, and a rectangular
pulse of force is applied at the right end during the first few instants.

Building the load as a table indexed by step number, as here, keeps the time
dependence explicit and readable. Note that the table must be filled for \emph{every} step the time loop will
read: a missing entry raises an error rather than returning zero. That is why
the loop opposite fills all \op{ndt}$+1$ entries, setting the two steps of
the pulse to \op{fimp} and all the others to zero.

\dindex{char[gement]}{loading history}%
For a long history it is more convenient to define an \op{evolution} and use
\op{char[gement]}, as in chapter~\ref{chap:plasticity}.
\end{castemcom}
\begin{castemlines}
\begin{verbatim}
AA = bloq 'DEPL' dgauche;

f0   = forc ddroite (0. 0.);
fimp = forc ddroite (-1. 0.);

F = table;
repeter bload (ndt + 1);
  k = &bload - 1;
  si ((k ega 1) ou (k ega 2));
    F . k = fimp;
  sino;
    F . k = f0;
  finsi;
fin bload;
\end{verbatim}
\end{castemlines}

\begin{castemcom}
The stiffness matrix $[K]$ and the mass matrix $[M]$ are assembled, and from
them the matrix $[S]=[M]+\beta\Delta t^{2}[K]$ -- the only one that will be
inverted during the algorithm, and the reason the implicit Newmark scheme
costs one factorisation rather than one per step.
\end{castemcom}
\begin{castemlines}
\begin{verbatim}
KK = rigi mod mat;
MM = mass mod mat;

SS = MM et
     ((beta * (dt ** 2.)) * KK);
\end{verbatim}
\end{castemlines}

\begin{castemcom}
Displacements $[U]$, velocities $[\dot{U}]$ and accelerations $[\ddot{U}]$ are
initialised. The initial acceleration is not chosen freely: it is obtained
from the equation of motion at $t=0$, i.e. by solving
$M\ddot{U}_0 = F_0 - K U_0$, which here reduces to $M\ddot{U}_0 = F_0$ since
$U_0=0$.
\end{castemcom}
\begin{castemlines}
\begin{verbatim}
U   = table;
dU  = table;
ddU = table;

U  . 0 = manu 'CHPO' dom 2
              'UX' 0. 'UY' 0.;
dU . 0 = manu 'CHPO' dom 2
              'UX' 0. 'UY' 0.;

ddU . 0 = reso (MM et AA) (F . 0);
\end{verbatim}
\end{castemlines}

\begin{castemcom}
And finally the time increment loop, in three stages: prediction, computation
of the acceleration, correction and update.

The three stages are written out so that the structure of the Newmark scheme
stays visible; in production \op{pasapas} does the same thing.
\end{castemcom}
\begin{castemlines}
\begin{verbatim}
repeter boucle ndt;
  n = &boucle - 1;

* ** prediction
  Upred  = (U . n)
    + (dt * (dU . n))
    + ((0.5 * dt * dt
        * (1. - (2.*beta))) * (ddU . n));
  dUpred = (dU . n)
    + ((dt * (1. - gamma)) * (ddU . n));

* ** acceleration
  ddU . (n + 1) = reso (SS et AA)
    ((F . (n+1)) - (KK * Upred));

* ** correction and update
  U  . (n + 1) = Upred
    + ((dt * dt * beta) * (ddU . (n+1)));
  dU . (n + 1) = dUpred
    + ((dt * gamma) * (ddU . (n+1)));
fin boucle;
\end{verbatim}
\end{castemlines}

\framepart{Summary}

\begin{gbox*}
\begin{itemize}
  \item \textbf{Adding time changes the loop, not the operators.} Both
  algorithms in this chapter have the same shape: assemble the constant
  matrices once outside the loop, keep the history in a \op{table} indexed
  by step number, and solve one linear system per step.
  \item \textbf{A constant time step is worth having.} The only matrix that
  is factorised is the one multiplying the unknown at $n+1$; keep the step
  constant and it is factorised once.
  \item \textbf{Two parameters decide whether the result means anything}:
  the scheme parameters -- $\theta$ for the heat equation,
  $(\beta,\gamma)$ for Newmark -- and the time step. $\theta\ge1/2$ and
  $\gamma=1/2$, $\beta=1/4$ are unconditionally stable and are the defaults
  to use until you have a reason not to.
  \item \textbf{Stability is not accuracy.} An over-large step gives a
  smooth, plausible and wrong answer.
\end{itemize}
\end{gbox*}

\framepart{Exercises}

\exercise{Stability of the $\theta$-scheme} Rerun the transient thermal
  computation with $\theta=0$ (explicit Euler) and increase $\Delta t$ until
  the solution oscillates. Compare the critical step with the theoretical
  stability limit. Then repeat with $\theta=1/2$ and $\theta=1$ and verify
  that no such limit exists.

\exercise{Energy conservation in Newmark} For the bar with no applied
  load after the initial pulse, compute at each step the kinetic energy
  $\tfrac12 \dot{U}^{T}M\dot{U}$ and the elastic energy
  $\tfrac12 U^{T}KU$, and plot their sum. Verify that it is conserved for
  $\gamma=1/2$ and dissipated for $\gamma>1/2$.

\exercise{Wave propagation} Plot the displacement of the free end of the
  bar against time and identify the round-trip travel time of the wave.
  Compare with $2L/c$.
